El Teorema~\ref{thm:limite} ya implica la positividad de $\lim_nf_h(n)$ para cada $h$ fijo.
En efecto, \eqref{eq:F} y las cotas del Teorema~\ref{thm:perfil}(a) y~(c) dan
$F(0)\ge\kappa^2>a_0/2$. La prueba directa del Teorema~\ref{thm:positiva} está en la
Sección~\ref{sec:positiva}. Lo enunciamos sólo para las dos últimas capas porque su papel es
mostrar la cuenta directa sin la cuenta inversa de la Sección~\ref{sec:inversa}; para todo $h$
fijo, el Teorema~\ref{thm:limite} da ya una conclusión más fuerte. La constante es conservadora y el
umbral $n_0$ es existencial.

Para entender esa prueba, volvamos al caso de la última reina. Si el cruce
$(r,c)$ está atacado, la inserción en esa casilla falla. Al retirar
una reina de $(i,j)$ podemos ocupar sus \emph{destinos} $(r,j)$ e $(i,c)$ si las dos reinas nuevas son
compatibles con el resto del tablero y entre sí.
Así se aprovechan tanto la fila $r$ como la columna $c$, aunque el cruce $(r,c)$
quede sin reina.

Para cualquier $h\ge1$, llamamos \emph{reparación de un movimiento} en $(r,c)$
a esta sustitución de una reina por dos. Una reparación legal lleva la colocación
$P\in\Om_h(n)$ a la capa $h-1$. La Figura~\ref{fig:reparacion-tablero}
muestra un caso en $4\times4$. El cruce vacío $(2,1)$ está atacado por la reina
de $(3,2)$. Retiramos la reina de $(0,0)$ y colocamos dos en $(2,0)$ y $(0,1)$.
El cruce original sigue bajo ataque y queda sin ocupar. La nueva colocación ya
tiene cuatro reinas y coincide con la Figura~\ref{fig:ejemplos}(a).

\begin{figure}[htbp]
\centering
\begin{tikzpicture}[tablero A]
 \begin{scope}
  \cuadriculaA{4}\trazarRejillaA{4}
  \ataqueA{3}{2}{2}{1}\aspaA{2}{1}
  \aroA{0}{0}\reinaA[colorAro]{0}{0}\reinaA{1}{3}\reinaA{3}{2}
  \panelA{2,4.38}{(a) Se retira una reina}
 \end{scope}
 \draw[-{Stealth[length=2mm]},tintaReina,line width=.9pt] (4.9,2)--(6.65,2);
 \begin{scope}[shift={(8,0)}]
  \cuadriculaA{4}\trazarRejillaA{4}
  \ataqueA{3}{2}{2}{1}\aspaA{2}{1}
  \reinaA{1}{3}\reinaA{3}{2}
  \aroA{0}{1}\reinaA[colorAro]{0}{1}
  \aroA{2}{0}\reinaA[colorAro]{2}{0}
  \panelA{2,4.38}{(b) Se colocan dos}
 \end{scope}
\end{tikzpicture}
\caption{Una reparación legal lleva de $\Omega_1(4)$ a $\Omega_0(4)$.
Los aros identifican las piezas que cambian. El aspa marca la misma casilla
$(2,1)$ en ambos dibujos. Leyendo de derecha a izquierda se recupera la operación
inversa con el par ordenado $(2,0),(0,1)$.}
\label{fig:reparacion-tablero}
\end{figure}


Llamaremos \emph{absorbedor} de ese cruce a una reina cuya sustitución
de esta manera es legal. En la Figura~\ref{fig:reparacion-tablero}(a), sólo
sirve la de $(0,0)$. Si retiramos la de $(1,3)$, uno de sus destinos es $(1,1)$.
Si retiramos la de $(3,2)$, uno de sus destinos es $(2,2)$. En ambos casos la
reina de $(0,0)$ permanece en su sitio y ataca ese destino. Tenemos, pues, un
cruce con cero inserciones directas y una reparación legal.

Sea $A^{rc}(P)$ el número de
reparaciones legales en $(r,c)$ y $d_{h,1}(P)=\sum A^{rc}(P)$, sumando sobre los $h^2$ cruces
vacíos; el subíndice $1$ alude a las reparaciones de un movimiento, en las que se retira una sola reina.
Escribimos $\bar r=(r+\frac12)/n$ y $\bar c=(c+\frac12)/n$ para los centros de las
casillas. Para expresar cómo varía el número de reparaciones con la posición del
cruce, usamos el perfil lleno. Denotamos su densidad por $\nu_1$: el subíndice
es la masa de reinas, igual a uno, que corresponde a $x=0$. Escribimos sus factores sin
el subíndice $x$ (la constante conserva el suyo, $b_0$) y ponemos
\begin{gather*}
 \nu_1=b_0\,r_R\,r_C\,u_D\,u_A,\qquad \psi(r,c)=u_D(r-c)\,u_A(r+c),\\
 a(r,c)=\int_{[0,1]^2}\nu_1(i,j)\,\psi(r,j)\,\psi(i,c)\,di\,dj .
\end{gather*}
Se tiene $\kappa^4\le a\le1$.

\begin{teorema}[Estructura típica]\label{thm:estructura}
Sea $h\ge1$ fijo y $P$ distribuida según $\mu_h$. Los factores $r_R,r_C$, la constante $b_0$ y
las funciones $\psi$ y $a$ son los del perfil lleno del Teorema~\ref{thm:perfil}, tomado
de~\cite{Perfil}.
\begin{enumerate}[(a)]
\item \emph{Reparaciones de un movimiento.} En probabilidad,
\[
 \max_{(r,c)\ \text{cruce vacío de }P}\ \Bigl|\frac{A^{rc}(P)}{n}-a(\bar r,\bar c)\Bigr|\longrightarrow0 .
\]
\item \emph{Posición media de los huecos.} Para toda función continua $\varphi$ en $[0,1]^2$,
\[
 \E_{\mu_h}\Bigl[\frac1{h^2}\sum_{(r,c)\ \text{cruce vacío}}\varphi(\bar r,\bar c)\Bigr]
 \longrightarrow\int_{[0,1]^2}\varphi(r,c)\,r_R(r)\,r_C(c)\,dr\,dc .
\]
Para $h=1$ esto es la convergencia débil de la ley de la posición del único cruce vacío.
\item \emph{Contador sin peso.} $\E_{\mu_h}[d_{h,1}]/(h^2n)\to B_1/b_0$, donde
$B_1=\int_{[0,1]^2}\psi$.
\end{enumerate}
\end{teorema}

\begin{observacion}[Alcance]\label{obs:alcance}
\begin{enumerate}[(i)]
\item Todos los enunciados son para $h$ y $K$ fijos cuando $n\to\infty$. No se obtienen tasas de
convergencia ni uniformidad para $K=K(n)\to\infty$, y no se determina la asintótica completa
de $Q(n)$.
\item En (b), para $h>1$ sólo se identifica la medida media de los cruces vacíos: no se afirma
que las filas y las columnas vacías sean asintóticamente independientes, ni que la medida empírica
de los cruces vacíos se concentre. En (c) se trata de
convergencia de medias: por (a), $d_{h,1}/n-\sum a(\bar r,\bar c)\to0$, y el segundo término
puede seguir siendo aleatorio, de modo que no se afirma concentración de $d_{h,1}/n$.
\item El valor $F(0)\approx0{,}0809$ procede de una resolución numérica determinista del
problema variacional~\cite[Sección~9]{Perfil}, sin certificación de error. A tamaño finito los valores
exactos, obtenidos por
enumeración exhaustiva, quedan lejos del límite (por ejemplo $f_1(15)\approx0{,}114$ y
$f_2(15)\approx0{,}153$;
Tabla~\ref{tab:ilustracion}); el teorema no da información sobre la velocidad de convergencia.
\item Los Teoremas~\ref{thm:limite}, \ref{thm:positiva} y~\ref{thm:estructura} usan el perfil del
Teorema~\ref{thm:perfil}, cuyas propiedades se toman de~\cite{Perfil}; el
Apéndice~\ref{app:perfil} resume las pruebas de las cotas y de la regularidad Lipschitz.
\end{enumerate}
\end{observacion}

\subsection{Idea de la demostración}\label{subsec:idea}

La Figura~\ref{fig:esquema} resume la estructura de la prueba.

Una misma colocación puede admitir varias reparaciones en un cruce. Contamos cada
operación legal una vez desde su colocación inicial y otra desde la final. En
la segunda descripción, partimos de $P'\in\Om_{h-1}$, elegimos el par ordenado
de reinas $(r,j),(i,c)$ y las sustituimos por $(i,j)$ para recuperar $P\in\Om_h$.
Asignamos al cruce $(r,c)$ el mismo peso $w(r,c)$ en ambas descripciones. La
doble cuenta de la Sección~\ref{sec:doblecuenta} da
\[
 N_h\,\E_{\mu_h}D_w=N_{h-1}\,\E_{\mu_{h-1}}I_w,\qquad\text{luego}\qquad
 f_h=\frac{n-h+1}{n}\,\frac{\E_{\mu_h}[D_w/(h^2n)]}{\E_{\mu_{h-1}}[I_w/n^2]},
\]
donde $D_w$ e $I_w$ son los contadores directo e inverso ponderados. Para usar el
cociente elegiremos un peso para el que el denominador tenga límite positivo.

Las dos medias se toman bajo leyes distintas: $\mu_h$ en la cuenta directa y
$\mu_{h-1}$ en la inversa. Para tratar ambas, la Sección~\ref{sec:referencia}
construye un proceso aleatorio para cada capa $h$. Coloca sucesivamente piezas reales
(reinas) y piezas virtuales, que no colocan ninguna reina pero desactivan las líneas que tocan,
hasta que no queda ninguna línea activa; los pesos proceden del perfil variacional. Denotemos por
$\nu'_h$ la ley de la colocación final en la capa $h$, condicionada al éxito del proceso. Una identidad
entrópica, las estimaciones
de los normalizadores (las sumas de pesos que normalizan cada elección) y la cota inferior de
primer orden de Simkin dan
$\KL{\mu_h}{\nu'_h}=o(n)$. En consecuencia, un evento con probabilidad a lo sumo
$e^{-cn}$ bajo $\nu'_h$, con $c>0$ fijo, tiene probabilidad $o(1)$ bajo $\mu_h$.
